% Options for packages loaded elsewhere
\PassOptionsToPackage{unicode}{hyperref}
\PassOptionsToPackage{hyphens}{url}
%
\documentclass[
]{article}
\usepackage{amsmath,amssymb}
\usepackage{iftex}
\ifPDFTeX
  \usepackage[T1]{fontenc}
  \usepackage[utf8]{inputenc}
  \usepackage{textcomp} % provide euro and other symbols
\else % if luatex or xetex
  \usepackage{unicode-math} % this also loads fontspec
  \defaultfontfeatures{Scale=MatchLowercase}
  \defaultfontfeatures[\rmfamily]{Ligatures=TeX,Scale=1}
\fi
\usepackage{lmodern}
\ifPDFTeX\else
  % xetex/luatex font selection
\fi
% Use upquote if available, for straight quotes in verbatim environments
\IfFileExists{upquote.sty}{\usepackage{upquote}}{}
\IfFileExists{microtype.sty}{% use microtype if available
  \usepackage[]{microtype}
  \UseMicrotypeSet[protrusion]{basicmath} % disable protrusion for tt fonts
}{}
\makeatletter
\@ifundefined{KOMAClassName}{% if non-KOMA class
  \IfFileExists{parskip.sty}{%
    \usepackage{parskip}
  }{% else
    \setlength{\parindent}{0pt}
    \setlength{\parskip}{6pt plus 2pt minus 1pt}}
}{% if KOMA class
  \KOMAoptions{parskip=half}}
\makeatother
\usepackage{xcolor}
\usepackage{longtable,booktabs,array}
\usepackage{calc} % for calculating minipage widths
% Correct order of tables after \paragraph or \subparagraph
\usepackage{etoolbox}
\makeatletter
\patchcmd\longtable{\par}{\if@noskipsec\mbox{}\fi\par}{}{}
\makeatother
% Allow footnotes in longtable head/foot
\IfFileExists{footnotehyper.sty}{\usepackage{footnotehyper}}{\usepackage{footnote}}
\makesavenoteenv{longtable}
\setlength{\emergencystretch}{3em} % prevent overfull lines
\providecommand{\tightlist}{%
  \setlength{\itemsep}{0pt}\setlength{\parskip}{0pt}}
\setcounter{secnumdepth}{-\maxdimen} % remove section numbering
\ifLuaTeX
  \usepackage{selnolig}  % disable illegal ligatures
\fi
\usepackage{bookmark}
\IfFileExists{xurl.sty}{\usepackage{xurl}}{} % add URL line breaks if available
\urlstyle{same}
\hypersetup{
  hidelinks,
  pdfcreator={LaTeX via pandoc}}

\author{}
\date{}

\begin{document}

\section{Cahier des Charges - Cyberspace IoT Security
Center}\label{cahier-des-charges---cyberspace-iot-security-center}

\textbf{Projet :} Prototype IoT de supervision, détection et réponse
cyber \textbf{Date :} 25 septembre 2026 \textbf{Technologies :}
Raspberry Pi (x2), MQTT (Mosquitto), Python, Flask, Docker, CI/CD

\begin{center}\rule{0.5\linewidth}{0.5pt}\end{center}

\subsection{1. Vision}\label{vision}

Cyberspace Horizon 2080 est un prototype de centre de sécurité
inter-systèmes. Il supervise des équipements IoT (capteurs, digicodes)
répartis dans un environnement spatial simulé. Il intercepte les signaux
MQTT, les chiffre, et les analyse en temps réel pour détecter et bloquer
les intrusions.

\textbf{Chaîne de valeur :} équipement → événement MQTT → détection IDS
→ alerte → réponse automatisée (confinement).

\textbf{Principes directeurs :} 1. \textbf{Zéro confiance envers le
transport.} Tout message MQTT est considéré hostile s'il n'est pas
chiffré (AES-128) et protégé par un jeton anti-rejeu. 2. \textbf{Réponse
Autonome.} L'isolement d'un équipement compromis est réalisé
instantanément par le serveur, sans intervention humaine requise. 3.
\textbf{Simplicité et robustesse.} Le système fonctionne intégralement
en mémoire vive locale (Edge Computing) pour garantir des temps de
réponse de l'ordre de la milliseconde.

\subsection{2. Contexte, hypothèses et
contraintes}\label{contexte-hypothuxe8ses-et-contraintes}

\begin{itemize}
\tightlist
\item
  Projet réalisé dans le cadre du Workshop Horizon 2080 : équipe de 4,
  durée courte.
\item
  \textbf{Hypothèse H1 :} Un conteneur Docker Mosquitto local sécurisé
  (identifiants stricts) est le cœur de la communication.
\item
  \textbf{Hypothèse H2 :} Le système est déployé sur deux cartes
  Raspberry Pi. La Board 1 (pi-center) agit comme Serveur Central. La
  Board 2 agit comme Client IoT (Digicode, Capteur PIR, Terminal de
  Chat).
\item
  \textbf{Contrainte :} Déploiement cible ARM64, réseau Wi-Fi local
  isolé.
\end{itemize}

\subsection{3. Objectifs et priorisation
(MoSCoW)}\label{objectifs-et-priorisation-moscow}

\begin{longtable}[]{@{}
  >{\raggedright\arraybackslash}p{(\columnwidth - 2\tabcolsep) * \real{0.4706}}
  >{\raggedright\arraybackslash}p{(\columnwidth - 2\tabcolsep) * \real{0.5294}}@{}}
\toprule\noalign{}
\begin{minipage}[b]{\linewidth}\raggedright
Niveau
\end{minipage} & \begin{minipage}[b]{\linewidth}\raggedright
Contenu
\end{minipage} \\
\midrule\noalign{}
\endhead
\bottomrule\noalign{}
\endlastfoot
\textbf{MUST} & Dashboard temps réel, Ingestion MQTT chiffrée, IDS
automatique, Détection de messages en clair, Détection de Replay
Attacks, Isolation automatique (Watchdog), CI/CD multi-architecture. \\
\textbf{SHOULD} & Capteur de mouvement physique réel (GPIO), Digicode
d'accès physique, Chat sécurisé inter-systèmes. \\
\textbf{WON'T} & Base de données persistante (SQLite abandonné au profit
de la RAM pour les performances), caméras virtuelles, authentification
par jetons web API. \\
\end{longtable}

\subsection{4. Architecture Globale}\label{architecture-globale}

\begin{itemize}
\tightlist
\item
  \textbf{Flux Sécurisé :} Capteur (PIR/Digicode) → Chiffrement AES-128
  Fernet + Horodatage → MQTT (QoS 1) → Serveur Central → IDS.
\item
  \textbf{Composants :}

  \begin{itemize}
  \tightlist
  \item
    \textbf{Serveur Central (Python/Flask) :} Orchestration, Dashboard
    Web.
  \item
    \textbf{Broker (Mosquitto) :} Transport asynchrone sécurisé
    (identifiants \texttt{v-client}).
  \item
    \textbf{IDS (Intrusion Detection System) :} Moteur de règles qui
    lève des alertes \texttt{CRITICAL} et isole les IPs/Devices
    compromis via le pare-feu logiciel.
  \item
    \textbf{Couche Physique (Raspberry Pi 2) :} Capteur Infrarouge (PIR)
    et script digicode interactif bloquant le sas physique.
  \end{itemize}
\end{itemize}

\subsection{5. Modèle de Menace \&
Contre-Mesures}\label{moduxe8le-de-menace-contre-mesures}

\begin{longtable}[]{@{}
  >{\raggedright\arraybackslash}p{(\columnwidth - 4\tabcolsep) * \real{0.1739}}
  >{\raggedright\arraybackslash}p{(\columnwidth - 4\tabcolsep) * \real{0.1739}}
  >{\raggedright\arraybackslash}p{(\columnwidth - 4\tabcolsep) * \real{0.6522}}@{}}
\toprule\noalign{}
\begin{minipage}[b]{\linewidth}\raggedright
Menace
\end{minipage} & \begin{minipage}[b]{\linewidth}\raggedright
Impact
\end{minipage} & \begin{minipage}[b]{\linewidth}\raggedright
Contre-mesure (Implémentée)
\end{minipage} \\
\midrule\noalign{}
\endhead
\bottomrule\noalign{}
\endlastfoot
\textbf{Message en clair injecté} & Contournement du système & L'IDS
détecte l'absence de chiffrement et lève une alerte
\texttt{UNENCRYPTED\_PAYLOAD}. \\
\textbf{Rejeu d'un message capturé} & Fausse commande d'ouverture & Le
payload inclut un \emph{Token} horodaté (TTL 300s). L'IDS le rejette
(\texttt{INVALID\_TOKEN}). \\
\textbf{Brute Force (Digicode)} & Accès non autorisé au SAS & Le script
bloque l'accès physique et remonte l'alerte \texttt{AUTH\_FAILURE} au
Dashboard. \\
\textbf{Intrusion Physique} & Sabotage des capteurs & Le capteur PIR
détecte un mouvement en zone verrouillée et alerte l'IDS. \\
\textbf{Comportement suspect} & Spam du réseau MQTT & Le script
\texttt{watchdog.py} place l'équipement en quarantaine totale
(\texttt{ISOLATED}). \\
\end{longtable}

\subsection{6. Critères de réussite de la Démonstration (5
minutes)}\label{crituxe8res-de-ruxe9ussite-de-la-duxe9monstration-5-minutes}

\begin{enumerate}
\def\labelenumi{\arabic{enumi}.}
\tightlist
\item
  \textbf{Trafic normal :} Les terminaux communiquent de manière
  chiffrée et apparaissent ``ONLINE'' sur le Dashboard.
\item
  \textbf{Capteur Physique :} Déclenchement d'un mouvement sur le
  capteur PIR sans taper le code \texttt{2080} déclenche une alerte au
  centre de contrôle.
\item
  \textbf{Attaque Simulée :} Lancement du script
  \texttt{test\_hacker.py}.
\item
  \textbf{Réaction de l'IDS :} L'alerte \texttt{CRITICAL} apparaît en
  temps réel sur l'interface Web (moins de 2 secondes).
\item
  \textbf{Isolation :} Le Dashboard confirme que l'équipement
  malveillant a été mis en statut \texttt{ISOLATED} et ne peut plus
  nuire.
\end{enumerate}

\end{document}
